% Generated by GrindEQ Word-to-LaTeX 
\documentclass{article} % use \documentstyle for old LaTeX compilers

\usepackage[utf8]{inputenc} % 'cp1252'-Western, 'cp1251'-Cyrillic, etc.
\usepackage[english]{babel} % 'french', 'german', 'spanish', 'danish', etc.
\usepackage{amsmath}
\usepackage{amssymb}
\usepackage{txfonts}
\usepackage{mathdots}
\usepackage[classicReIm]{kpfonts}
\usepackage{graphicx}

%%% remove comment delimiter ('%') and select graphics package
%%% for DVI output:
\usepackage[dvips]{graphicx}
%%% or for PDF output:
%\usepackage[pdftex]{graphicx}
%%% or for old LaTeX compilers:
%\usepackage[dvips]{graphics}


\begin{document}

%\selectlanguage{english} % remove comment delimiter ('%') and select language if required


\noindent 

\noindent 

\noindent 

\noindent 

\noindent \textbf{Design and mechanical property research of composite cylindrical lattice structures with excellent stiffness and strength inspired by plant stems}\footnote{$\mathrm{\ast}$Corresponding author. E-mail address: jx@hit.edu.cn (J. Xiong).}

\noindent Xun Wang${}^{a,b}$, Jian Xiong${}^{a,b}$${}^{\mathrm{\ast }}$

\noindent ${}^{a}$ National Key Laboratory of Science and Technology for Advanced Composites in Special Environments, Center for Composite Materials and Structures, Harbin Institute of Technology, Harbin, 150080, PR China

\noindent ${}^{b}$ Structural Lightweighting and Composites Laboratory (StruCM Lab), Center for Composite Materials and Structures, Harbin Institute of Technology, Harbin, 150080, PR China

\noindent \textbf{Highlight}

\eqref{GrindEQ__1_} A truss cylindrical lattice structure inspired by plant stalks is proposed.

\eqref{GrindEQ__2_} Theoretical prediction models for its compressive modulus and compressive strength are established, and the validity of the models is verified via experiments and numerical simulations.

\eqref{GrindEQ__3_} The strength, stiffness and impact energy absorption characteristics of this type of truss cylindrical lattice structure under low-density conditions are revealed.

\noindent 
\subsection{Abstract}

Lightweight lattice structures have attracted considerable research attention due to their high geometric symmetry and superior lightweight performance. Inspired by the excellent stiffness of plant stalks, four cylindrical lattice structures with distinct topological configurations are designed to meet the demand for load-bearing capacity and impact resistance of curved aerospace structures. Theoretical prediction models for compressive modulus and compressive strength are established based on the energy method. The load-bearing and energy absorption characteristics of the cylindrical lattice structures under quasi-static compression and impact loads are analyzed via finite element simulations. Composite truss cylindrical lattice structures are fabricated using selective laser sintering technology, and quasi-static compression and low-velocity impact experiments are carried out to verify the accuracy of the theoretical models and numerical simulations. The results demonstrate that relative density and topological configuration exert significant effects on the stiffness, strength, and impact energy absorption properties of truss cylindrical lattice structures. This work clarifies the load-bearing and energy absorption characteristics of truss cylindrical lattice structures under low-density conditions, and provides a theoretical foundation for the lightweight design and load-bearing performance optimization of curved aerospace structures.

\noindent Keywords: Lattice structures; Compressive modulus; Compressive strength; Impact energy absorption

\noindent 
\subsection{1  Introduction}

Lightweight lattice structures are of great research significance in fields including aerospace and rail transportation due to their lightweight nature, high specific strength, high specific stiffness and designability${}^{[1-10]}$.Truss lattice structures are a class of porous materials composed of periodically arranged strut units interconnected at nodes.The design of such structures is based on a tension-dominated deformation mechanism: when subjected to external loads, the constituent struts primarily undergo axial tensile or compressive deformation rather than bending deformation, thus achieving high stiffness and strength at low relative densities. Accordingly, truss lattice structures have become one of the most extensively investigated lattice structures in recent years${}^{[11]}$.

In recent decades, extensive studies have been carried out on the mechanical properties of truss lattice structures with various configurations${}^{[12-19]}$. For truss lattice structures, experimental, numerical simulation and theoretical methods are commonly employed to investigate their mechanical performance under quasi-static and dynamic compression conditions${}^{[20-29]}$. In terms of quasi-static mechanical properties, Li et al.${}^{[30]}$ enhanced the mechanical performance of lattice structures by introducing a cross-strut strategy and established a prediction model for the maximum compressive strength. Chen et al.${}^{[31]}$ reconstructed simple cubic, face-centered cubic (FCC) and body-centered cubic (BCC) lattices into self-similar nested configurations, and clarified the effect of the nesting ratio on compressive response and energy absorption performance. Han et al.${}^{[32]}$ verified that BCC-derived lattice structures have superior load-bearing capacity and energy absorption characteristics compared with the original BCC structure. Kothari et al.${}^{[33]}$ adopted a design combining I-shaped struts and node reinforcement, which significantly improved the specific stiffness, specific yield strength and energy absorption capacity of the lattice. The curved BCC structure proposed by Sasagawa et al.${}^{[34]}$ exhibits far greater flexibility and a wider tunable range of mechanical properties than both conventional BCC and octet truss lattices at the same relative density. For BCC lattice structures under compression with large deformation strain, Chen et al.${}^{[35]}$ derived the correlations between strut slenderness ratio and key mechanical properties including Young's modulus and yield strength. A systematic study via finite element modeling was conducted to explore the equivalent mechanical properties of BCC lattice structures under compression. The novel hollow octet truss lattice structure developed by Zhang et al.${}^{[36]}$ addresses the instability issue of conventional tension-dominated lattices at low relative densities, and its remarkable energy absorption performance has been validated through experiments and numerical simulations. Lan et al.${}^{[37]}$ achieved the synergy of high energy absorption capacity and stable deformation behavior based on a multi-dimensional gradient design of octet truss topological lattice structures. The vertex-modified BCC lattice structure proposed by Wang et al.${}^{[38]}$ outperforms conventional truss lattice structures in deformation stability, energy absorption capacity and strength.

Research on dynamic mechanical properties focuses on the energy absorption characteristics of lightweight lattice structures under impact loads, offering valuable insights for the design of lightweight lattice structures subjected to impact loading.The vertex-offset face-centered cubic (FCC) lattice structure proposed by Wang et al. ${}^{[39]}$ exhibits enhanced energy absorption capacity and a stable stress-strain response under axial impact, as well as excellent local deformation resistance under oblique impact, achieving efficient load transfer. Its bionic design strategy provides a new approach for the optimization of multi-directional mechanical properties. Dejean et al. ${}^{[40]}$ analyzed the uniaxial quasi-static compression and dynamic impact properties of octet truss lattice structures using a combined experimental and numerical method, and revealed the effect of loading rates on their mechanical performance.Sun et al. ${}^{[41]}$ obtained the dynamic and quasi-static compression response curves of octahedral, rhombic dodecahedral and novel hybrid lattice structures via drop-weight tests and universal compression testing machines, respectively, and clarified the regulation mechanism of deformation rate on the stress-strain response and energy absorption performance of the structures. Chen et al. ${}^{[42]}$ analyzed the effects of Kelvin and auxetic lattice structures on cushioning behavior and load distribution through impact experiments and finite element simulations. Zimmerman et al. ${}^{[43]}$ investigated the yield characteristics of octet truss lattice structures under impact loads and discussed the influence of face sheet thickness on the load attenuation mechanism, providing references for structural parameter optimization under impact loading.As for truss cylindrical lattice structures, current research mainly focuses on their mechanical properties under quasi-static compression, while studies on their energy absorption characteristics under dynamic impact loads remain limited.

Regarding cylindrical lattice structures, Gao et al. ${}^{[44]}$ validated the accuracy of the theoretical model via numerical simulations and experiments based on a novel negative Poisson's ratio cylindrical double-arrow honeycomb structure. Zhang et al. ${}^{[45]}$ investigated the peak stress, energy absorption performance and load efficiency of hollow cylindrical and cubic lattice structures by adjusting the twist angle. Tian et al. ${}^{[46]}$ developed a bio-inspired dual-phase cylindrical lattice, which achieves higher specific energy absorption than the corresponding single-phase structure. Pan et al. ${}^{[47]}$ proposed a reinforced cylindrical negative stiffness structure with internally filled inclined beams, and analyzed its static and dynamic mechanical properties. Li et al. ${}^{[48]}$ constructed a cylindrical lattice using rhombic dodecahedron unit cells, and acquired its force-displacement response and energy absorption characteristics through impact tests. Sood et al. ${}^{[49]}$ analyzed the mechanical properties and energy absorption features of lattice cylindrical structures via static and dynamic tests combined with numerical simulations. Gong et al. ${}^{[50]}$ integrated three types of concave cylindrical lattice structures with distinct Poisson's ratio properties with triply periodic minimal surfaces, and characterized their mechanical properties through quasi-static compression tests. The results demonstrate that the hybrid design yields a remarkable improvement in energy absorption performance. As for truss cylindrical lattice structures, existing studies mainly focus on finite element analysis and experimental investigations, while systematic research on their theoretical models remains relatively limited.

Based on the above literature review, existing studies on lightweight lattice structures are predominantly focused on regular cubic configurations, which have attracted widespread research interest owing to their high geometric symmetry and superior lightweight performance. Furthermore, regular cubic truss lattice structures fail to meet the impact resistance and energy absorption requirements of cylindrical curved components in the aerospace field under practical service conditions. Nevertheless, there remain notable research gaps in the field of composite truss cylindrical lattice structures. Specifically, the theoretical models and impact resistance characteristics of truss cylindrical lattice structures under low-density conditions have not been thoroughly investigated. Accordingly, systematic research on the load-bearing mechanisms and deformation laws of composite truss cylindrical lattice structures can provide a theoretical foundation for lightweight design in the aerospace sector. Critical issues to be resolved include realizing the design evolution of truss lattice configurations from regular cubic forms to cylindrical curved surfaces, and constructing theoretical models for the compressive strength and stiffness of truss cylindrical lattice structures at low densities. To address the aforementioned research gaps, this paper first investigates low-density composite cylindrical lattice structures inspired by plant stalks. The relative densities of various truss cylindrical lattice structures are calculated, and theoretical models for the compressive modulus and compressive strength of composite truss cylindrical lattice structures are established in combination with structural geometric characteristics. Finite element models of composite truss cylindrical lattice structures are developed to study the mechanical response of the structures under quasi-static compression and impact loads. Subsequently, quasi-static compression tests and impact tests are carried out to acquire the load-displacement curves of the structures. Finally, comparative analyses are performed among theoretical calculation results, numerical simulation results and experimental data. The mechanical properties and deformation mechanisms of the structures under low-density conditions are revealed, the mechanical performance of truss cylindrical lattice structures with different relative densities is clarified, and the validity of the proposed theoretical models is verified.

\noindent 
\subsection{2  Design process of truss lattice structures}

\noindent 
\paragraph{2.1 Design process from cubic configuration to curved configuration of truss lattice structures}

To realize the precise design of traditional truss lattice configurations from regular cubic forms to cylindrical curved surfaces, a low-density truss cylindrical lattice structure is proposed. Figure 1 illustrates the design process of transforming truss lattice structures from cubic configurations to cylindrical curved surface configurations. Through geometric adaptation and structural reconstruction, the transition to cylindrical curved surface components is achieved, so as to meet the demands for impact resistance and energy absorption of curved components in aerospace and other fields.

First, typical cubic truss lattice units are presented, which feature distinct geometric symmetry and periodic arrangement. Afterwards, a series of geometric transformation and mapping operations are applied to gradually unfold the cubic units and attach them to the cylindrical curved surface. During this process, the connectivity of truss members and the basic topological characteristics of node distribution are retained. Meanwhile, the length of truss members and the position of nodes are adjusted to form the truss cylindrical lattice structure. By properly controlling member diameter, node spacing and interlayer arrangement, the resulting cylindrical lattice structure maintains excellent deformation resistance even at a low density. This systematic design method from cubic configurations to cylindrical curved surfaces provides a research framework for subsequent mechanical modeling, simulation analysis and experimental verification of truss cylindrical lattice structures.

\noindent \includegraphics*[width=4.51in, height=4.05in]{image1}

Fig. 2-1  The design process from truss lattice cubic configuration to curved surface

\noindent 
\paragraph{2.2 Configuration of composite truss cylindrical lattice structures}

In nature, plants such as wheat, dandelion and horsetail can sustain their own weight, thanks to the excellent load-bearing capacity and deformation resistance of their stems. As shown in Figure 2(a), staggered internal struts form periodic supporting frames. The truss cylindrical lattice structure is adopted to achieve uniform load transfer and enhance the overall deformation resistance. Accordingly, drawing on the superior stiffness characteristics of plant stems, this paper designs several bionic composite truss cylindrical lattice structures, which deliver excellent load-bearing performance and stiffness. Figure 2(b) presents optical micrographs that demonstrate the cylindrical lattice morphology of plant stems. At a low relative density, these truss cylindrical lattice structures fully leverage the advantages of light weight and high stiffness, and thus possess outstanding mechanical properties. This research provides effective ideas for the lightweight design of curved-surface structures in relevant engineering fields.

\noindent \includegraphics*[width=3.54in, height=2.70in]{image2}

Fig. 2  The cylindrical porous structures in nature

\noindent \includegraphics*[width=4.37in, height=4.69in]{image3}

Fig. 3  The design process and geometric parameters of composite truss cylindrical lattice structures

The design process and basic dimensional parameters of composite truss cylindrical lattice structures are illustrated in Figure 3. The proposed structures are classified and named as body-centered cubic cylindrical lattice structure, face-centered cubic cylindrical lattice structure, octahedral cylindrical lattice structure and tesseract cylindrical lattice structure respectively.

\noindent 
\paragraph{2.3 Relative density of truss cylindrical lattice structures}

The basic geometric parameters of composite truss cylindrical lattice structures cover the height of the unit cell and the number of structural layers, which are denoted by $l_{i} $ and \textit{n} respectively with \textit{n}=6 . They also include the radii of the outer circle and inner circle, represented by $R_{i} $ and $r_{i} $ correspondingly. The truss members of the above cylindrical lattice structures adopt circular cross-sections, and the diameter of each member is defined as $d_{i} $. Here, \textit{i}=1, 2, 3, 4 stand for body-centered cubic, face-centered cubic, octahedral and tesseract cylindrical lattice structures in sequence. The relationships between the main geometric parameters are as follows:
\begin{equation} \label{GrindEQ__1_} 
A_{i0} =\pi \left(R_{i}^{2} -r_{i}^{2} \right) 
\end{equation} 
\begin{equation} \label{GrindEQ__2_} 
A_{i} =\frac{1}{4} \pi d_{i}^{2}  
\end{equation} 
\begin{equation} \label{GrindEQ__3_} 
h_{i} =R_{i} -r_{i}  
\end{equation} 

For the body-centered cubic cylindrical lattice structure, its relative density ($\bar{\rho }_{1} $) can be expressed as:
\begin{equation} \label{GrindEQ__4_} 
\bar{\rho }_{1} =\frac{A_{1} \left[2\left(n+1\right)\pi \left(R_{1} +r_{1} \right)+12\left(n+1\right)h_{1} +24nl_{1} +12\left(N_{1,1} l_{1,1} +N_{1,2} l_{1,2} \right)\right]-V_{c1} }{nA_{10} l_{1} }  
\end{equation} 
where, in Eq.\eqref{GrindEQ__4_}, $l_{1,1} $, $l_{1,2} $ and $V_{c1} $ denote the total lengths of all straight struts and the overlapping volume of the body-centered cubic cylindrical lattice structure, respectively. Their respective expressions are given as follows:
\[\begin{array}{l} {N_{1,1} =4n\begin{array}{cc} {,} & {l_{1,1} =\frac{1}{2} } \end{array}\left[R_{1}^{2} +\left(5-\sqrt{2} -\sqrt{6} \right)r_{1}^{2} +\left(2-\sqrt{2} -\sqrt{6} \right)R_{1} r_{1} +l_{1}^{2} \right]^{{\raise0.7ex\hbox{$ 1 $}\!\mathord{\left/{\vphantom{1 2}}\right.\kern-\nulldelimiterspace}\!\lower0.7ex\hbox{$ 2 $}} } ,} \\ {N_{1,2} =4n\begin{array}{cc} {,} & {l_{1,2} =\frac{1}{2} } \end{array}\left[\left(5-\sqrt{2} -\sqrt{6} \right)R_{1}^{2} +r_{1}^{2} +\left(2-\sqrt{2} -\sqrt{6} \right)R_{1} r_{1} +l_{1}^{2} \right]^{{\raise0.7ex\hbox{$ 1 $}\!\mathord{\left/{\vphantom{1 2}}\right.\kern-\nulldelimiterspace}\!\lower0.7ex\hbox{$ 2 $}} } ,} \\ {V_{c1} =12\left(2\sqrt{6} n+\frac{4}{3} n+\frac{2}{3} \right)d_{1}^{3} .} \end{array}\] 

For the face-centered cubic cylindrical lattice structure, its relative density ($\bar{\rho }_{2} $) can be expressed as:
\begin{equation} \label{GrindEQ__5_} 
\bar{\rho }_{2} =\frac{A_{2} \left[2\left(n+1\right)\pi \left(R_{2} +r_{2} \right)+12\left(n+1\right)h_{2} +24nl_{2} +12\sum _{j=1}^{4}N_{2,j}  l_{2,j} \right]-V_{c2} }{nA_{20} l_{2} }  
\end{equation} 
where, in Eq. \eqref{GrindEQ__5_}, $l_{2,i} $ (\textit{i}=1, 2, 3, 4) and $V_{c2} $ denote the total length of all straight struts and the overlapping volume of the face-centered cubic cylindrical lattice structure, respectively. Their respective expressions are given as follows:
\[\sum _{j=1}^{4}N_{2,j}  l_{2,j} =N_{2,1} l_{2,1} +N_{2,2} l_{2,2} +N_{2,3} l_{2,3} +N_{2,4} l_{2,4} \begin{array}{cc} {,} & {N_{2,1} } \end{array}=2n\begin{array}{cc} {,} & {l_{2,1} =} \end{array}\left[\left(2-\sqrt{3} \right)r_{2}^{2} +l_{2}^{2} \right]^{{\raise0.7ex\hbox{$ 1 $}\!\mathord{\left/{\vphantom{1 2}}\right.\kern-\nulldelimiterspace}\!\lower0.7ex\hbox{$ 2 $}} } ,N_{2,2} =2n\begin{array}{cc} {,} & {l_{2,2} =} \end{array}\left[\left(2-\sqrt{3} \right)R_{2}^{2} +l_{2}^{2} \right]^{{\raise0.7ex\hbox{$ 1 $}\!\mathord{\left/{\vphantom{1 2}}\right.\kern-\nulldelimiterspace}\!\lower0.7ex\hbox{$ 2 $}} } \begin{array}{cc} {,} & {N_{2,3} } \end{array}=2\left(n+1\right)\begin{array}{cc} {,} & {l_{2,3} } \end{array}=\left(R_{2}^{2} -\sqrt{3} R_{2} r_{2} +r_{2}^{2} \right)^{{\raise0.7ex\hbox{$ 1 $}\!\mathord{\left/{\vphantom{1 2}}\right.\kern-\nulldelimiterspace}\!\lower0.7ex\hbox{$ 2 $}} } ,\] 
\[N_{2,4} =2n\begin{array}{cc} {,} & {l_{2,4} =} \end{array}\left(R_{2}^{2} -2R_{2} r_{2} +r_{2}^{2} +l_{2}^{2} \right)^{{\raise0.7ex\hbox{$ 1 $}\!\mathord{\left/{\vphantom{1 2}}\right.\kern-\nulldelimiterspace}\!\lower0.7ex\hbox{$ 2 $}} } \begin{array}{cc} {,} & {V_{c2} } \end{array}=12\left(3\sqrt{6} n+4n+\frac{4}{3} \right)d_{2}^{3} .\] 

For the octahedral cylindrical lattice structure, its relative density ($\bar{\rho }_{3} $) can be expressed as:
\begin{equation} \label{GrindEQ__6_} 
\bar{\rho }_{3} =\frac{12\cdot \left(A_{3} \sum _{j=1}^{9}N_{3,j} l_{3,j}  \right)-V_{c3} }{nA_{30} l_{3} }  
\end{equation} 
where, in Eq. \eqref{GrindEQ__6_}, $l_{3,i} $ ($i=1,2,\cdot \cdot \cdot ,9$) and $V_{c3} $ denote the total length of all straight struts and the overlapping volume of the octahedral cylindrical lattice structure, respectively. Their respective expressions are given as follows:
\[\begin{array}{l} {\sum _{j=1}^{9}N_{3,j} l_{3,j}  =N_{3,1} l_{3,1} +N_{3,2} l_{3,2} +N_{3,3} l_{3,3} +N_{3,4} l_{3,4} +N_{3,5} l_{3,5} +N_{3,6} l_{3,6} +N_{3,7} l_{3,7} +N_{3,8} l_{3,8} +N_{3,9} l_{3,9} } \\ {N_{3,1} =2n\begin{array}{cc} {,} & {l_{3,1} =} \end{array}\left[\left(2-\sqrt{3} \right)r_{3}^{2} +l_{3}^{2} \right]^{{\raise0.7ex\hbox{$ 1 $}\!\mathord{\left/{\vphantom{1 2}}\right.\kern-\nulldelimiterspace}\!\lower0.7ex\hbox{$ 2 $}} } \begin{array}{cc} {,} & {N_{3,2} } \end{array}=2n\begin{array}{cc} {,} & {l_{3,2} =} \end{array}\left[\left(2-\sqrt{3} \right)R_{3}^{2} +l_{3}^{2} \right]^{{\raise0.7ex\hbox{$ 1 $}\!\mathord{\left/{\vphantom{1 2}}\right.\kern-\nulldelimiterspace}\!\lower0.7ex\hbox{$ 2 $}} } ,} \end{array}N_{3,3} =2\left(n+1\right)\begin{array}{cc} {,} & {l_{3,3} =} \end{array}\left(R_{3}^{2} -\sqrt{3} R_{3} r_{3} +r_{3}^{2} \right)^{{\raise0.7ex\hbox{$ 1 $}\!\mathord{\left/{\vphantom{1 2}}\right.\kern-\nulldelimiterspace}\!\lower0.7ex\hbox{$ 2 $}} } \begin{array}{cc} {,} & {N_{3,4} } \end{array}=2n\begin{array}{cc} {,} & {l_{3,4} =} \end{array}\left(R_{3}^{2} -2R_{3} r_{3} +r_{3}^{2} +l_{3}^{2} \right)^{{\raise0.7ex\hbox{$ 1 $}\!\mathord{\left/{\vphantom{1 2}}\right.\kern-\nulldelimiterspace}\!\lower0.7ex\hbox{$ 2 $}} } ,N_{3,5} =2n\begin{array}{cc} {,} & {l_{3,5} =} \end{array}\frac{1}{2} \left[\left(2+\sqrt{3} \right)\left(\frac{R_{3} -r_{3} }{R_{3} +r_{3} } \right)r_{3}^{2} +l_{3}^{2} \right]^{{\raise0.7ex\hbox{$ 1 $}\!\mathord{\left/{\vphantom{1 2}}\right.\kern-\nulldelimiterspace}\!\lower0.7ex\hbox{$ 2 $}} } ,N_{3,6} =2n\begin{array}{cc} {,} & {l_{3,6} =} \end{array}\frac{1}{2} \left[\left(2+\sqrt{3} \right)\left(\frac{R_{3} -r_{3} }{R_{3} +r_{3} } \right)R_{3}^{2} +l_{3}^{2} \right]^{{\raise0.7ex\hbox{$ 1 $}\!\mathord{\left/{\vphantom{1 2}}\right.\kern-\nulldelimiterspace}\!\lower0.7ex\hbox{$ 2 $}} } ,\begin{array}{l} {N_{3,9} =4n\begin{array}{cc} {,} & {l_{3,9} =} \end{array}\frac{1}{2} \left(\frac{R_{3}^{4} -2\sqrt{3} R_{3}^{3} r_{3} +6R_{3}^{2} r_{3}^{2} -2\sqrt{3} R_{3} r_{3}^{3} +r_{3}^{4} }{R_{3}^{2} +2R_{3} r_{3} +r_{3}^{2} } +l_{3}^{2} \right)^{{\raise0.7ex\hbox{$ 1 $}\!\mathord{\left/{\vphantom{1 2}}\right.\kern-\nulldelimiterspace}\!\lower0.7ex\hbox{$ 2 $}} } ,} \\ {V_{c3} =12\left(3\sqrt{6} n+\frac{8}{3} n+\frac{2}{3} \right)d_{3}^{3} .} \end{array}\] 

For the tesseract cylindrical lattice structure, its relative density ($\bar{\rho }_{4} $) can be expressed as:
\begin{equation} \label{GrindEQ__7_} 
\bar{\rho }_{4} =\frac{A_{4} \left[2\left(n+1\right)\pi \left(R_{4} +r_{4} \right)+12\left(n+1\right)h_{4} +24nl_{4} +12\left(\sum _{j=1}^{6}N_{4,j}  l_{4,j} \right)\right]-V_{c4} }{nA_{40} l_{4} }  
\end{equation} 
where, in Eq. \eqref{GrindEQ__7_}, $l_{4,i} $ (\textit{i}=1, 2, 3, 4 ) and $V_{c4} $ denote the total length of all straight struts and the overlapping volume of the tesseract cylindrical lattice structure, respectively. Their respective expressions are given as follows:
\[\begin{array}{l} {\sum _{j=1}^{6}N_{4,j}  l_{4,1} +N_{4,2} l_{4,2} +N_{4,3} l_{4,3} +N_{4,4} l_{4,4} +N_{4,5} l_{4,5} +N_{4,6} l_{4,6} } \\ {N_{4,1} =4n\begin{array}{cc} {,} & {l_{4,1} =\frac{1}{4} } \end{array}\left[R_{4}^{2} +\left(5-\sqrt{2} -\sqrt{6} \right)r_{4}^{2} +\left(2-\sqrt{2} -\sqrt{6} \right)R_{4} r_{4} +l_{4}^{2} \right]^{{\raise0.7ex\hbox{$ 1 $}\!\mathord{\left/{\vphantom{1 2}}\right.\kern-\nulldelimiterspace}\!\lower0.7ex\hbox{$ 2 $}} } ,} \end{array}\] 
\[\begin{array}{l} {N_{4,2} =4n\begin{array}{cc} {,} & {l_{4,2} =} \end{array}\frac{1}{4} \left[\left(5-\sqrt{2} -\sqrt{6} \right)R_{4}^{2} +r_{4}^{2} +\left(2-\sqrt{2} -\sqrt{6} \right)R_{4} r_{4} +l_{4}^{2} \right]^{{\raise0.7ex\hbox{$ 1 $}\!\mathord{\left/{\vphantom{1 2}}\right.\kern-\nulldelimiterspace}\!\lower0.7ex\hbox{$ 2 $}} } ,} \\ {N_{4,3} =4n\begin{array}{cc} {,} & {l_{4,3} =} \end{array}\frac{1}{2} \left(R_{4} -r_{4} \right)\begin{array}{cc} {,} & {N_{4,4} } \end{array}=4n\begin{array}{cc} {,} & {l_{4,4} =} \end{array}\frac{1}{2} l_{4} \begin{array}{cc} {,} & {N_{4,5} =2n\begin{array}{cc} {,} & {l_{4,5} =} \end{array}\frac{\pi }{12} r_{4} ,} \end{array}} \end{array}\] 
\[N_{4,6} =2n\begin{array}{cc} {,} & {l_{4,6} =} \end{array}\frac{\pi }{12} R_{4} \begin{array}{cc} {,} & {V_{c4} } \end{array}=12\left(16\sqrt{6} n+\frac{10}{3} n+\frac{2}{3} \right)d_{4}^{3} .\] 

Among the four composite truss cylindrical lattice structures, the outer radius $R_{i} $ , inner radius $r_{i} $ and structural height $l_{i} $ (\textit{i}=1, 2, 3, 4) are fixed at 38 mm, 27 mm and 15 mm, respectively. Other geometric parameters are adjusted according to the relative density, as shown in Table 1. The relative density is determined by two independent variables, i.e., the unit cell height $l_{i} $ and strut diameter $d_{i} $. For a constant relative density, modifying either variable can regulate the mechanical properties of the composite truss cylindrical lattice structures.

Table 1  Dimensions and relative density of composite truss cylindrical lattice structures

\begin{tabular}{|p{0.7in}|p{0.6in}|p{0.6in}|p{0.6in}|p{0.6in}|p{0.7in}|} \hline 
Cylindrical lattice structures & $R_{i} $\newline /(mm) & $r_{i} $\newline /(mm) & $l_{i} $\newline /(mm) & $d_{i} $\newline /(mm) & Relative density $\bar{\rho }$ \\ \hline 
Body-centered & 38 & 22 & 15 & 2 & 0.130 \\ \hline 
 & 42.5 & 22.5 & 20 & 2 & 0.092 \\ \hline 
 & 52.5 & 22.5 & 25 & 2 & 0.057 \\ \hline 
Face-centered & 38 & 22 & 15 & 2 & 0.156 \\ \hline 
 & 42.5 & 22.5 & 20 & 2 & 0.113 \\ \hline 
 & 52.5 & 22.5 & 25 & 2 & 0.072 \\ \hline 
Octahedral & 38 & 22 & 15 & 2 & 0.229 \\ \hline 
 & 42.5 & 22.5 & 20 & 2 & 0.160 \\ \hline 
 & 52.5 & 22.5 & 25 & 2 & 0.0999 \\ \hline 
Tesseract & 38 & 22 & 15 & 2 & 0.091 \\ \hline 
 & 42.5 & 22.5 & 20 & 2 & 0.073 \\ \hline 
 & 52.5 & 22.5 & 25 & 2 & 0.051 \\ \hline 
\end{tabular}

\includegraphics*[width=2.74in, height=2.10in]{image4}\includegraphics*[width=2.74in, height=2.10in]{image5}

\noindent \includegraphics*[width=2.74in, height=2.10in]{image6}\includegraphics*[width=2.74in, height=2.10in]{image7}

Fig. 4  The relative density of the composite truss cylindrical attice structures as a function of geometric parameters: (a) body-centered; (b) face-centered; (c) octahedral; (d) tesseract.

Figure 4 presents the relationship between relative density and geometric dimensions of composite truss cylindrical lattice structures. It can be seen from the figure that the relative density is governed by a variety of geometric parameters, such as unit cell height and strut diameter. Taking the four categories of truss cylindrical lattice structures, namely body-centered cubic, face-centered cubic, octahedral and tesseract structures, as research objects, their outer radius, inner radius and overall height are fixed at 38 mm, 27 mm and 15 mm respectively. Under such conditions, the relative density changes along with the variation of strut diameter.

When the outer circle diameter increases within a certain range, the relative density shows an overall rising trend. This phenomenon demonstrates that a larger outer diameter corresponds to a denser mass distribution of the structure, and the material filling level is higher under the same overall external dimensions. The variation curves of relative density differ from one structural type to another. For example, the tesseract cylindrical lattice structure possesses the minimum relative density among all the structures. It is demonstrated that different structural types impose disparate effects on relative density. Owing to their distinctive geometric compositions, each structure displays different evolutionary features of relative density when the same geometric parameters are adjusted. Consequently, the relative density of the structure can be regulated by adjusting geometric parameters, so as to further enhance the mechanical properties of the structure.

\noindent 
\paragraph{2.4 Theoretical model of composite truss cylindrical lattice structures}

\noindent 
\subparagraph{2.4.1 Compressive modulus}

In this section, theoretical models for compressive modulus and compressive strength are established based on the energy method. For the body-centered cubic cylindrical lattice structure, the area moment of inertia of the circular strut is given as:
\begin{equation} \label{GrindEQ__8_} 
I_{1} =\frac{1}{64} \pi d_{1}^{4}  
\end{equation} 

As shown in Figure 5(a), the force $F_{a} $ acting on the axial strut can be expressed as:
\begin{equation} \label{GrindEQ__9_} 
F_{a} =\frac{\delta _{1} E_{s} A_{10} }{l_{1} }  
\end{equation} 
where $\delta _{1} $ represents the compressive deformation.

The forces $F_{b}^{p} $ and $F_{c}^{p} $ acting on the inclined struts are expressed as:
\begin{equation} \label{GrindEQ__10_} 
\left\{\begin{array}{l} {F_{b}^{p} =\frac{3E_{s} I_{1} \delta _{1} l_{1} }{4\kappa _{1,1}^{4} } } \\ {F_{c}^{p} =\frac{3E_{s} I_{1} \delta _{1} l_{1} }{4\kappa _{1,2}^{4} } } \end{array}\right.  
\end{equation} 

The forces $F_{b}^{v} $ and $F_{c}^{v} $ perpendicular to the inclined struts are expressed as:
\begin{equation} \label{GrindEQ__11_} 
\left\{\begin{array}{l} {F_{b}^{v} =\frac{3E_{s} I_{1} l_{1,1} \delta _{1} }{2\kappa _{1,1}^{4} } } \\ {F_{c}^{v} =\frac{3E_{s} I_{1} l_{1,2} \delta _{1} }{2\kappa _{1,2}^{4} } } \end{array}\right.  
\end{equation} 
where the lengths lengths $\kappa _{1,1} $ and $\kappa _{1,2} $ of the inclined struts A${}_{1}$D${}_{1}$ and C${}_{1}$D${}_{1}$ are expressed as:
\[\begin{array}{l} {\kappa _{1,1} =\frac{1}{2} \left[R_{1}^{2} +\left(5-\sqrt{2} -\sqrt{6} \right)r_{1}^{2} +\left(2-\sqrt{2} -\sqrt{6} \right)R_{1} r_{1} \right]^{{\raise0.7ex\hbox{$ 1 $}\!\mathord{\left/{\vphantom{1 2}}\right.\kern-\nulldelimiterspace}\!\lower0.7ex\hbox{$ 2 $}} } } \\ {\kappa _{1,2} =\frac{1}{2} \left[\left(5-\sqrt{2} -\sqrt{6} \right)R_{1}^{2} +r_{1}^{2} +\left(2-\sqrt{2} -\sqrt{6} \right)R_{1} r_{1} \right]^{{\raise0.7ex\hbox{$ 1 $}\!\mathord{\left/{\vphantom{1 2}}\right.\kern-\nulldelimiterspace}\!\lower0.7ex\hbox{$ 2 $}} } } \end{array}\] 

The strain energy $U_{a} $ of the axial strut is expressed as:
\begin{equation} \label{GrindEQ__12_} 
U_{a} =\frac{F_{AB}^{2} l_{1} }{2A_{10} E_{s} }  
\end{equation} 
\includegraphics*[width=5.90in, height=1.49in]{image8}

\noindent \includegraphics*[width=5.87in, height=1.49in]{image9}

\noindent \includegraphics*[width=5.84in, height=2.64in]{image10}

\noindent \includegraphics*[width=5.82in, height=1.50in]{image11}

Fig. 5  A schematic diagram of the force analysis of the composite truss cylindrical lattice structures

The axial strain energies $U_{b}^{p} $ and $U_{c}^{p} $ along the inclined struts are expressed as:
\begin{equation} \label{GrindEQ__13_} 
\left\{\begin{array}{l} {U_{b}^{p} =\frac{F_{b}^{p} {}^{2} \kappa _{1,1} }{2A_{10} E_{s} } } \\ {U_{c}^{p} =\frac{F_{c}^{p} {}^{2} \kappa _{1,2} }{2A_{10} E_{s} } } \end{array}\right.  
\end{equation} 

The bending strain energies $U_{b}^{v} $ and $U_{c}^{v} $ perpendicular to the inclined struts are expressed as:
\begin{equation} \label{GrindEQ__14_} 
\left\{\begin{array}{l} {U_{b}^{v} =\frac{F_{b}^{v} {}^{2} \kappa _{1,1}^{3} }{6E_{s} I_{1} } } \\ {U_{c}^{v} =\frac{F_{c}^{v} {}^{2} \kappa _{1,2}^{3} }{6E_{s} I_{1} } } \end{array}\right.  
\end{equation} 

The total strain energy $U_{1} $ of the body-centered cubic cylindrical lattice structure is expressed as:
\begin{equation} \label{GrindEQ__15_} 
U_{1} =24\cdot U_{a} +48\left(U_{b}^{p} +U_{b}^{v} +U_{c}^{p} +U_{c}^{v} \right) 
\end{equation} 

The equivalent strain energy $U_{eq1} $ of the body-centered cubic cylindrical lattice structure is expressed as:
\begin{equation} \label{GrindEQ__16_} 
U_{eq1} =\frac{1}{2} F_{a} \delta _{1} =\frac{\delta _{1}^{2} E_{1}^{*} A_{10} }{2l_{1} }  
\end{equation} 

Based on the energy method, the total strain energy $U_{1} $ of the body-centered cubic cylindrical lattice structure is equal to its equivalent strain energy $U_{eq1} $, as expressed below:
\begin{equation} \label{GrindEQ__17_} 
U_{1} =U_{eq1}  
\end{equation} 

Substituting Eqs. \eqref{GrindEQ__15_} and \eqref{GrindEQ__16_} into Eq. \eqref{GrindEQ__17_}, we obtain the expression for the compressive modulus $E_{1}^{*} $ of the body-centered cubic cylindrical lattice structure as follows:
\begin{equation} \label{GrindEQ__18_} 
E_{1}^{*} =\left[\frac{6d_{1}^{2} }{R_{1}^{2} -r_{1}^{2} } +\frac{9d_{1}^{4} l_{1,1}^{2} l_{1} }{16\left(R_{1}^{2} -r_{1}^{2} \right)\kappa _{1,1}^{5} } +\frac{9d_{1}^{4} l_{1,2}^{2} l_{1} }{16\left(R_{1}^{2} -r_{1}^{2} \right)\kappa _{1,2}^{5} } +\frac{27d_{1}^{6} l_{1}^{3} }{1024\left(R_{1}^{2} -r_{1}^{2} \right)\kappa _{1,1}^{7} } +\frac{27d_{1}^{6} l_{1}^{3} }{1024\left(R_{1}^{2} -r_{1}^{2} \right)\kappa _{1,2}^{7} } \right]\cdot E_{s}  
\end{equation} 

As shown in Figure 5(b), consistent with the body-centered cubic cylindrical lattice structure, the compressive modulus $E_{2}^{*} $ of the face-centered cubic cylindrical lattice structure is obtained as:
\begin{equation} \label{GrindEQ__19_} 
\begin{array}{l} {E_{2}^{*} =\left[\frac{6d_{2}^{2} }{R_{2}^{2} -r_{2}^{2} } +\frac{9d_{2}^{4} \left(R_{2} -r_{2} \right)^{2} l_{2} }{2l_{2,4}^{5} \left(R_{2}^{2} -r_{2}^{2} \right)} +\frac{27d_{2}^{6} l_{2}^{3} }{8l_{2,4}^{7} \left(R_{2}^{2} -r_{2}^{2} \right)} +\frac{18d_{2}^{4} \kappa _{2,1}^{2} l_{2} }{l_{2,1}^{5} \left(R_{2}^{2} -r_{2}^{2} \right)} \right. } \\ {+\frac{18d_{2}^{4} \kappa _{2,2}^{2} l_{2} }{l_{2,2}^{5} \left(R_{2}^{2} -r_{2}^{2} \right)} +\frac{6d_{2}^{2} l_{2}^{3} }{l_{2,1}^{3} \left(R_{2}^{2} -r_{2}^{2} \right)} +\left. \frac{8d_{2}^{2} l_{2}^{3} }{l_{2,2}^{3} \left(R_{2}^{2} -r_{2}^{2} \right)} \right]\cdot E_{s} } \end{array} 
\end{equation} 
where parameters $\kappa _{2,1} $ and $\kappa _{2,2} $ in Eq. \eqref{GrindEQ__19_} are expressed as:
\[\kappa _{2,1} =\frac{\sqrt{6} -\sqrt{2} }{4} r_{2} \begin{array}{cc} {,} & {\kappa _{2,2} } \end{array}=\frac{\sqrt{6} -\sqrt{2} }{4} R_{2} \] 

As shown in Figure 5(c), the corresponding compressive modulus $E_{3}^{*} $ of the octahedral cylindrical lattice structure is expressed as:
\begin{equation} \label{GrindEQ__20_} 
\begin{array}{l} {E_{3}^{*} =\left[\frac{9d_{3}^{4} \left(R_{3} -r_{3} \right)^{2} l_{3} }{64l_{3,4}^{5} \left(R_{3}^{2} -r_{3}^{2} \right)} +\right. \frac{27d_{3}^{6} l_{3}^{3} }{1024l_{3,4}^{7} \left(R_{3}^{2} -r_{3}^{2} \right)} +\frac{9d_{3}^{4} \kappa _{4,1}^{2} l_{3} }{16l_{3,1}^{5} \left(R_{3}^{2} -r_{3}^{2} \right)} +\frac{9d_{3}^{4} \kappa _{4,2}^{2} l_{3} }{16l_{3,2}^{5} \left(R_{3}^{2} -r_{3}^{2} \right)} } \\ {+\frac{3d_{3}^{2} l_{3}^{3} }{4l_{3,1}^{3} \left(R_{3}^{2} -r_{3}^{2} \right)} +\frac{3d_{3}^{2} l_{3}^{3} }{4l_{3,2}^{3} \left(R_{3}^{2} -r_{3}^{2} \right)} +\frac{3d_{3}^{2} l_{3}^{3} }{8\kappa _{5,1}^{3} \left(R_{3}^{2} -r_{3}^{2} \right)} +\frac{3d_{3}^{2} l_{3}^{3} }{8\kappa _{5,2}^{3} \left(R_{3}^{2} -r_{3}^{2} \right)} } \\ {+\frac{9d_{3}^{4} \kappa _{6,1}^{2} l_{3} }{32\kappa _{5,1}^{5} \left(R_{3}^{2} -r_{3}^{2} \right)} +\frac{9d_{3}^{4} \kappa _{6,2}^{2} l_{3} }{32\kappa _{5,2}^{5} \left(R_{3}^{2} -r_{3}^{2} \right)} +\frac{9d_{3}^{4} \kappa _{3,2}^{2} l_{3} }{16\kappa _{3,1}^{5} \left(R_{3}^{2} -r_{3}^{2} \right)} +\left. \frac{27d_{3}^{6} l_{3}^{3} }{1024\kappa _{3,1}^{7} \left(R_{3}^{2} -r_{3}^{2} \right)} \right]\cdot E_{s} } \end{array} 
\end{equation} 
where the parameters in Eq. \eqref{GrindEQ__20_} are expressed as:
\[\kappa _{3,1} =\frac{1}{2} \left(\frac{R_{3}^{4} -2\sqrt{3} R_{3}^{3} r_{3} +6R_{3}^{2} r_{3}^{2} -2\sqrt{3} R_{3} r_{3}^{3} +r_{3}^{4} }{R_{3}^{2} +2R_{3} r_{3} +r_{3}^{2} } +l_{3}^{2} \right)^{{\raise0.7ex\hbox{$ 1 $}\!\mathord{\left/{\vphantom{1 2}}\right.\kern-\nulldelimiterspace}\!\lower0.7ex\hbox{$ 2 $}} } \begin{array}{cc} {,} & {\kappa _{4,1} } \end{array}=\frac{\sqrt{6} -\sqrt{2} }{4} r_{3} ,\] 
\[\kappa _{3,2} =\frac{1}{2\left(R_{3} +r_{3} \right)} \left(R_{3}^{4} -2\sqrt{3} R_{3}^{3} r_{3} +6R_{3}^{2} r_{3}^{2} -2\sqrt{3} R_{3} r_{3}^{3} +r_{3}^{4} \right)^{{\raise0.7ex\hbox{$ 1 $}\!\mathord{\left/{\vphantom{1 2}}\right.\kern-\nulldelimiterspace}\!\lower0.7ex\hbox{$ 2 $}} } \begin{array}{cc} {,} & {\kappa _{4,2} } \end{array}=\frac{\sqrt{6} -\sqrt{2} }{4} R_{3} .\] 
\[\begin{array}{l} {\kappa _{5,1} =\frac{1}{2} \left[\left(2+\sqrt{3} \right)\left(\frac{R_{3} -r_{3} }{R_{3} +r_{3} } \right)^{2} r_{3}^{2} +l_{3}^{2} \right]^{{\raise0.7ex\hbox{$ 1 $}\!\mathord{\left/{\vphantom{1 2}}\right.\kern-\nulldelimiterspace}\!\lower0.7ex\hbox{$ 2 $}} } \begin{array}{cc} {,} & {\kappa _{5,2} } \end{array}=\frac{1}{2} \left[\left(2+\sqrt{3} \right)\left(\frac{R_{3} -r_{3} }{R_{3} +r_{3} } \right)^{2} R_{3}^{2} +l_{3}^{2} \right]^{{\raise0.7ex\hbox{$ 1 $}\!\mathord{\left/{\vphantom{1 2}}\right.\kern-\nulldelimiterspace}\!\lower0.7ex\hbox{$ 2 $}} } } \\ {\kappa _{6,1} =\frac{\left(\sqrt{6} +\sqrt{2} \right)\left(R_{3} -r_{3} \right)}{4\left(R_{3} +r_{3} \right)} r_{3} \begin{array}{cc} {,} & {\kappa _{6,2} } \end{array}=\frac{\left(\sqrt{6} +\sqrt{2} \right)\left(R_{3} -r_{3} \right)}{4\left(R_{3} +r_{3} \right)} R_{3} } \end{array}\] 

As shown in Figure 5(d), the expression for the compressive modulus $E_{4}^{*} $ of the tesseract cylindrical lattice structure is presented below:
\begin{equation} \label{GrindEQ__21_} 
\begin{array}{l} {E_{4}^{*} =\left[\frac{6d_{4}^{2} }{R_{4}^{2} -r_{4}^{2} } +\frac{27d_{4}^{4} \kappa _{7,1}^{2} l_{4}^{2} }{256l_{4,1}^{6} \left(R_{4}^{2} -r_{4}^{2} \right)} +\frac{27d_{4}^{4} \kappa _{7,2}^{2} l_{4}^{2} }{256l_{4,2}^{6} \left(R_{4}^{2} -r_{4}^{2} \right)} +\frac{9d_{4}^{4} \kappa _{7,1}^{2} l_{4} }{64l_{4,1}^{5} \left(R_{4}^{2} -r_{4}^{2} \right)} \right. } \\ {+\frac{9d_{4}^{4} \kappa _{7,2}^{2} l_{4} }{64l_{4,2}^{5} \left(R_{4}^{2} -r_{4}^{2} \right)} +\frac{27d_{4}^{6} l_{4}^{3} }{16384l_{4,1}^{7} \left(R_{4}^{2} -r_{4}^{2} \right)} +\left. \frac{27d_{4}^{6} l_{4}^{3} }{16384l_{4,2}^{7} \left(R_{4}^{2} -r_{4}^{2} \right)} \right]\cdot E_{s} } \end{array} 
\end{equation} 
where the lengths $\kappa _{7,1} $ and $\kappa _{7,2} $ of the inclined struts A${}_{4}$D${}_{4}$ and E${}_{4}$F${}_{4}$ in Equation \eqref{GrindEQ__21_} are expressed as:
\[\begin{array}{l} {\kappa _{7,1} =\frac{1}{4} \left[\left(5-\sqrt{2} -\sqrt{6} \right)r_{4}^{2} +R_{4}^{2} +\left(2-\sqrt{2} -\sqrt{6} \right)R_{4} r_{4} \right]^{{\raise0.7ex\hbox{$ 1 $}\!\mathord{\left/{\vphantom{1 2}}\right.\kern-\nulldelimiterspace}\!\lower0.7ex\hbox{$ 2 $}} } ,} \\ {\kappa _{7,2} =\frac{1}{4} \left[\left(5-\sqrt{2} -\sqrt{6} \right)R_{4}^{2} +r_{4}^{2} +\left(2-\sqrt{2} -\sqrt{6} \right)R_{4} r_{4} \right]^{{\raise0.7ex\hbox{$ 1 $}\!\mathord{\left/{\vphantom{1 2}}\right.\kern-\nulldelimiterspace}\!\lower0.7ex\hbox{$ 2 $}} } } \end{array}\] 

\subparagraph{2.4.2 Compressive strength}

From Eq. \eqref{GrindEQ__18_}, the stress $\sigma _{1} $ of the body-centered cubic cylindrical lattice structure is expressed as:
\begin{equation} \label{GrindEQ__22_} 
\sigma _{1} =\frac{6d_{1}^{2} \delta _{1} E_{s} }{\left(R_{1}^{2} -r_{1}^{2} \right)l_{1} } +\frac{9d_{1}^{4} l_{1,1}^{2} \delta _{1} E_{s} }{16\left(R_{1}^{2} -r_{1}^{2} \right)\kappa _{1,1}^{5} } +\frac{9d_{1}^{4} l_{1,2}^{2} \delta _{1} E_{s} }{16\left(R_{1}^{2} -r_{1}^{2} \right)\kappa _{1,2}^{5} } +\frac{27d_{1}^{6} l_{1}^{2} \delta _{1} E_{s} }{1024\left(R_{1}^{2} -r_{1}^{2} \right)\kappa _{1,1}^{7} } +\frac{27d_{1}^{6} l_{1}^{2} \delta _{1} E_{s} }{1024\left(R_{1}^{2} -r_{1}^{2} \right)\kappa _{1,2}^{7} }  
\end{equation} 

As shown in Figure 4(a), the stress $\sigma _{a} $ of the axial strut is expressed as:
\begin{equation} \label{GrindEQ__23_} 
\sigma _{a} =\frac{\delta _{1} E_{s} }{l_{1} }  
\end{equation} 

The stresses $\sigma _{b} $ and $\sigma _{c} $ of the inclined struts are given as:
\begin{equation} \label{GrindEQ__24_} 
\left\{\begin{array}{l} {\sigma _{b} =\frac{3E_{s} l_{1,1} \delta _{1} }{\kappa _{1,1}^{3} } +\frac{3E_{s} d_{1}^{2} \delta _{1} l_{1} }{64\kappa _{1,1}^{4} } } \\ {\sigma _{c} =\frac{3E_{s} l_{1,2} \delta _{1} }{\kappa _{1,2}^{3} } +\frac{3E_{s} d_{1}^{2} \delta _{1} l_{1} }{64\kappa _{1,2}^{4} } } \end{array}\right.  
\end{equation} 

Substituting Eqs. \eqref{GrindEQ__23_} and \eqref{GrindEQ__24_} into Eq. \eqref{GrindEQ__22_}, the stress $\sigma _{1} $ of the body-centered cubic cylindrical lattice structure is derived as:
\begin{equation} \label{GrindEQ__25_} 
\sigma _{1} =\frac{3d_{1}^{2} }{R_{1}^{2} -r_{1}^{2} } \left[2\sigma _{a} +\frac{64d_{1} \kappa _{1,1}^{2} l_{1,1}^{2} +3d_{1}^{3} l_{1}^{2} }{16\kappa _{1,1}^{3} \left(d_{1} l_{1} +16\kappa _{1,1} l_{1,1} \right)} \sigma _{b} +\frac{64d_{1} \kappa _{1,2}^{2} l_{1,2}^{2} +3d_{1}^{3} l_{1}^{2} }{16\kappa _{1,2}^{3} \left(d_{1} l_{1} +16\kappa _{1,2} l_{1,2} \right)} \sigma _{c} \right] 
\end{equation} 

Using the mixed criterion for the failure strengths of three struts, the failure strength $\sigma _{1}^{*} $ of the body-centered cubic cylindrical lattice structure is obtained and expressed as:
\begin{equation} \label{GrindEQ__26_} 
\sigma _{1}^{*} =\frac{3d_{1}^{2} }{R_{1}^{2} -r_{1}^{2} } \left[2\cdot \sigma _{a1}^{pk} +\frac{64d_{1} \kappa _{11}^{2} l_{11}^{2} +3d_{1}^{3} l_{1}^{2} }{16\kappa _{11}^{3} \left(d_{1} l_{1} +16\kappa _{1,1} l_{1,1} \right)} \sigma _{b1}^{pk} +\frac{64d_{1} \kappa _{12}^{2} l_{12}^{2} +3d_{1}^{3} l_{1}^{2} }{16\kappa _{12}^{3} \left(d_{1} l_{1} +16\kappa _{1,2} l_{1,2} \right)} \sigma _{c1}^{pk} \right] 
\end{equation} 
where the plastic buckling stresses $\sigma _{a1}^{pk} $, $\sigma _{b1}^{pk} $ and $\sigma _{c1}^{pk} $ of the three types of struts are expressed as:
\[\sigma _{a1}^{pk} =\frac{25\pi ^{2} d_{1}^{2} }{196l_{1}^{2} } E_{s} \begin{array}{cc} {,} & {\sigma _{b1}^{pk} } \end{array}=\frac{25\pi ^{2} d_{1}^{2} }{196\kappa _{1,1}^{2} } E_{s} \begin{array}{cc} {,} & {\sigma _{c1}^{pk} } \end{array}=\frac{25\pi ^{2} d_{1}^{2} }{196\kappa _{1,2}^{2} } E_{s} \] 

Similar to the body-centered cubic cylindrical lattice structure, as shown in Figure 5(b), the failure strength $\sigma _{2}^{*} $ of the face-centered cubic cylindrical lattice structure can be derived using the mixed criterion for the failure strengths of the four struts, which is expressed as:
\begin{equation} \label{GrindEQ__27_} 
\begin{array}{l} {\sigma _{2}^{*} =\frac{3d_{2}^{2} }{R_{2}^{2} -r_{2}^{2} } \left\{2\right. \sigma _{a2}^{pk} +\frac{16d_{2} l_{2,4}^{2} \left(R_{2} -r_{2} \right)^{2} +3d_{2}^{3} l_{2}^{2} }{2l_{2,4}^{3} \left[d_{2} l_{2} +8l_{2,4} \left(R_{2} -r_{2} \right)\right]} \sigma _{b2}^{pk} } \\ {+\frac{3d_{2}^{2} \kappa _{2,1}^{2} +l_{2}^{2} l_{2,1}^{2} }{l_{2,1}^{2} \left(3\kappa _{2,1} d_{2} +0.5l_{2} l_{2,1} \right)} \sigma _{c2}^{pk} +\frac{3d_{2}^{2} \kappa _{2,2}^{2} +l_{2}^{2} l_{2,2}^{2} }{l_{2,2}^{2} \left(3\kappa _{2,2} d_{2} +0.5l_{2} l_{2,2} \right)} \left. \sigma _{d2}^{pk} \right\}} \end{array} 
\end{equation} 
where the plastic buckling stresses $\sigma _{a2}^{pk} $, $\sigma _{b2}^{pk} $, $\sigma _{c2}^{pk} $ and $\sigma _{d2}^{pk} $ of the four types of struts are expressed as:
\[\sigma _{a2}^{pk} =\frac{25\pi ^{2} d_{2}^{2} }{196l_{2}^{2} } E_{s} \begin{array}{cc} {,} & {\sigma _{b2}^{pk} } \end{array}=\frac{25\pi ^{2} d_{2}^{2} }{49l_{2,4}^{2} } E_{s} \begin{array}{cc} {,} & {\sigma _{c2}^{pk} } \end{array}=\frac{\pi ^{2} d_{2}^{2} }{4l_{2,1}^{2} } E_{s} \begin{array}{cc} {,} & {\sigma _{d2}^{pk} } \end{array}=\frac{\pi ^{2} d_{2}^{2} }{4l_{2,2}^{2} } E_{s} \] 

As shown in Figure 5(c), the failure strength $\sigma _{3}^{*} $ of the octahedral cylindrical lattice structure can be derived using the mixed criterion for the failure strengths of the six struts, which is expressed as:
\begin{equation} \label{GrindEQ__28_} 
\begin{array}{l} {\sigma _{3}^{*} =\frac{3d_{2}^{2} }{4\left(R_{2}^{2} -r_{2}^{2} \right)} \left[\frac{64d_{3} l_{34}^{2} \left(R_{4} -r_{4} \right)^{2} +3d_{3}^{3} l_{3}^{2} }{4l_{34}^{3} \left[d_{3} l_{3} +16l_{3,4} \left(R_{4} -r_{4} \right)\right]} \sigma _{a3}^{pk} \right. +\frac{3d_{3}^{2} \kappa _{41}^{2} +4l_{3}^{2} l_{31}^{2} }{\kappa _{71}^{2} \left(3\kappa _{4,1} d_{3} +l_{3} l_{3,1} \right)} \sigma _{b3}^{pk} } \\ {+\frac{3d_{3}^{2} \kappa _{42}^{2} +4l_{3}^{2} l_{32}^{2} }{l_{32}^{2} \left(3\kappa _{4,2} d_{3} +l_{3} l_{3,2} \right)} \sigma _{c3}^{pk} +\frac{3d_{3}^{2} \kappa _{61}^{2} +4l_{3}^{2} \kappa _{51}^{2} }{2\kappa _{51}^{2} \left(3\kappa _{6,1} d_{3} +l_{3} \kappa _{5,1} \right)} \sigma _{d3}^{pk} } \\ {+\frac{3d_{3}^{2} \kappa _{62}^{2} +4l_{3}^{2} \kappa _{52}^{2} }{2\kappa _{52}^{2} \left(3\kappa _{6,2} d_{3} +l_{3} \kappa _{5,2} \right)} \sigma _{e3}^{pk} +\frac{64d_{3} \kappa _{31}^{2} \kappa _{32}^{2} +3d_{3}^{3} l_{3}^{2} }{4\kappa _{31}^{3} \left(d_{3} l_{3} +16\kappa _{3,1} \kappa _{3,2} \right)} \left. \sigma _{f3}^{pk} \right]} \end{array} 
\end{equation} 
where the plastic buckling stresses $\sigma _{a3}^{pk} $, $\sigma _{b3}^{pk} $, $\sigma _{c3}^{pk} $ $\sigma _{d3}^{pk} $, $\sigma _{e3}^{pk} $,and $\sigma _{f3}^{pk} $ of the six types of struts are expressed as:
\[\sigma _{a3}^{pk} =\frac{25\pi ^{2} d_{3}^{2} }{196l_{3,4}^{2} } E_{s} \begin{array}{cc} {,} & {\sigma _{b3}^{pk} } \end{array}=\frac{\pi ^{2} d_{3}^{2} }{16l_{3,1}^{2} } E_{s} \begin{array}{cc} {,} & {\sigma _{c3}^{pk} } \end{array}=\frac{\pi ^{2} d_{3}^{2} }{16l_{3,2}^{2} } E_{s} ,\] 
\[\sigma _{d3}^{pk} =\frac{\pi ^{2} d_{3}^{2} }{16\kappa _{5,1}^{2} } E_{s} \begin{array}{cc} {,} & {\sigma _{e3}^{pk} } \end{array}=\frac{\pi ^{2} d_{3}^{2} }{16\kappa _{5,2}^{2} } E_{s} \begin{array}{cc} {,} & {\sigma _{f3}^{pk} } \end{array}=\frac{25\pi ^{2} d_{3}^{2} }{196\kappa _{3,1}^{2} } E_{s} \] 

As shown in Figure 5(d), the failure strength $\sigma _{4}^{*} $ of the tesseract cylindrical lattice structure can be derived using the mixed criterion for the failure strengths of the four struts, which is expressed as:
\begin{equation} \label{GrindEQ__29_} 
\begin{array}{l} {\sigma _{4}^{*} =\frac{3d_{4}^{2} }{R_{4}^{2} -r_{4}^{2} } \left[2\sigma _{a4}^{pk} \right. +\frac{256d_{4} l_{41}^{2} \kappa _{71}^{2} +3d_{4}^{3} l_{4}^{2} }{64l_{41}^{3} \left(d_{4} l_{4} +32l_{4,1} \kappa _{7,1} \right)} \sigma _{b4}^{pk} } \\ {+\frac{256d_{4} l_{42}^{2} \kappa _{72}^{2} +3d_{4}^{3} l_{4}^{2} }{64l_{42}^{3} \left(d_{4} l_{4} +32l_{4,2} \kappa _{7,2} \right)} \sigma _{c4}^{pk} +\frac{3d_{4} l_{4} \left(l_{41}^{3} \kappa _{7,2} +\kappa _{7,1} l_{42}^{3} \right)}{32l_{41}^{3} l_{42}^{3} } \left. \sigma _{d4}^{pk} \right]} \end{array} 
\end{equation} 
where the plastic buckling stresses $\sigma _{a4}^{pk} $, $\sigma _{b4}^{pk} $, $\sigma _{c4}^{pk} $ and $\sigma _{d4}^{pk} $ of the four types of struts are expressed as:
\[\sigma _{a4}^{pk} =\frac{25\pi ^{2} d_{4}^{2} }{196l_{4}^{2} } E_{s} \begin{array}{cc} {,} & {\sigma _{b4}^{pk} } \end{array}=\frac{25\pi ^{2} d_{4}^{2} }{196l_{4,1}^{2} } E_{s} \begin{array}{cc} {,} & {\sigma _{c4}^{pk} } \end{array}=\frac{25\pi ^{2} d_{4}^{2} }{196l_{4,2}^{2} } E_{s} \begin{array}{cc} {,} & {\sigma _{d4}^{pk} } \end{array}=\frac{25\pi ^{2} d_{4}^{2} }{49l_{4}^{2} } E_{s} \] 

\subsection{3 Finite element model of composite truss cylindrical lattice structures}

To further explore the mechanical properties of composite truss cylindrical lattice structures, verify the stiffness and strength theoretical models proposed in previous sections, and compare the predictions with experimental results, finite element analysis is adopted in this section to conduct numerical simulations on several types of such structures. First, the geometric models of the truss cylindrical lattice structures are established strictly according to the dimensions of actual specimens, and the detailed dimensional parameters are listed in Table 1. Subsequently, mesh generation is performed on the finite element models, where the tetrahedral solid element C3D10M is used throughout this study.

\noindent \includegraphics*[width=3.68in, height=3.12in]{image12}

Fig. 6  Mesh convergence analysis of the finite element model

The element models of three types of low-density truss cylindrical lattice structures were established using the analysis module ABAQUS/Explicit to determine their elastic modulus and strength under different low densities. Figure 6 presents the convergence analysis results of mesh size on the compressive strength in finite element simulation. To improve the efficiency of mesh convergence analysis for composite truss cylindrical lattice structures, the body-centered cubic cylindrical lattice structure was selected as the research object with geometric dimensions of $R_{1} $=38mm, $r_{1} $=22mm, $l_{1} $=15mm and $d_{1} $=2mm. The simulation results reveal that variations in mesh element size exert a slight influence on the final analysis results. Accordingly, a mesh size of 1 mm was adopted, which can satisfy the accuracy requirements of finite element simulation while ensuring overall computational efficiency.

\noindent \includegraphics*[width=2.85in, height=2.25in]{image13}\includegraphics*[width=2.83in, height=2.25in]{image14}

\noindent \includegraphics*[width=2.81in, height=2.33in]{image15}\includegraphics*[width=2.83in, height=2.33in]{image16}

Fig. 7  The boundary conditions and mesh configurations of the composite truss cylindrical lattice structures: (a) body-centered; (b) face-centered; (c) octahedral; (d) tesseract.

An elastoplastic model was adopted to characterize the stiffness and strength of truss cylindrical lattice structures under quasi-static compression, and the Dynamic, Explicit analysis step was employed in the finite element simulation. For boundary conditions, the lattice specimens were placed between two rigid compression platens, which were defined as discrete rigid bodies in the model. Corresponding degrees of freedom and boundary constraints were assigned to their reference points (RPs). To replicate uniform compression loading, the lower platen was fully fixed, while the upper platen moved downward along the \textit{y}-direction, as illustrated in Figure 7. General contact was defined between the platens and lattice structures. Hard contact was specified for normal behaviour, and the penalty function method was used for tangential behaviour with a friction coefficient of 0.15. The simulated elastic modulus of composite truss cylindrical lattice structures was calculated from the slope of the stress-strain curve within the elastic deformation stage.

To investigate the impact resistance of composite truss cylindrical lattice structures, finite element models were established to analyze their deformation and energy absorption characteristics under impact loading. The explicit dynamic module of ABAQUS was utilized for numerical modelling to compare the impact performance of different composite truss cylindrical lattice structures. The Johnson-Cook (J-C) plasticity model was adopted in the finite element analysis to describe the strain rate hardening effect of materials under impact loading. Since the thermal effect was not taken into account in the J-C constitutive model for the present dynamic impact simulation, the yield stress of the material under dynamic impact can be simplified to the following expression:
\begin{equation} \label{GrindEQ__30_} 
\sigma _{c} =\left[a+b\varepsilon _{p}^{n} \right]\cdot \left[1+c\ln \left(\frac{\dot{\varepsilon }_{p} }{\dot{\varepsilon }} \right)\right] 
\end{equation} 

In this equation, $a$ is the yield strength of the material, while $b$ and $n$ represent the strain rate hardening constant and strain rate hardening exponent, respectively. All these parameters are determined by fitting the stress-strain curves obtained from tensile tests on standard specimens in accordance with ASTM standards. In addition, $\dot{\varepsilon }_{p} $ and $\dot{\varepsilon }$ denote the dynamic strain rate and reference test strain rate, and $\varepsilon _{p} $ is the cumulative strain rate. According to the experimental results reported by Mars ${}^{[}$${}^{50}$${}^{]}$, the strain rate sensitivity constant $c$ is set to 0.002. The contact properties and mesh generation remain identical to those used in the quasi-static compression simulation. For the interaction settings under impact loading, a mass of 7.73 kg is assigned to the reference point (RP) of the upper steel plate, and a downward velocity load is applied to implement impact loading on the truss cylindrical lattice structures and extract stress and strain data during finite element simulation. The lower steel plate is constrained in the same manner as in the quasi-static compression tests.

\noindent 
\subsection{4 Experimental methods for composite truss cylindrical lattice structures}

Dog-bone specimens and composite truss cylindrical lattice structures with various configurations were fabricated via 3D printing using nylon PA6 reinforced with short glass fibers. Tensile tests were conducted on standard dog-bone specimens to establish the constitutive relation of the matrix material, which provides support for subsequent numerical simulations. Quasi-static compression and low-velocity impact tests were performed to validate the simulation results, and the experimental data of different composite truss cylindrical lattice structures were acquired.

\noindent 
\paragraph{4.1 Fabrication of composite truss cylindrical lattice structures}

The composite truss cylindrical lattice structures were fabricated from short glass fiber-reinforced nylon PA6 using selective laser sintering with a laser power of 200 W. After manufacturing, all specimens were sandblasted for 20 minutes to remove residual nylon PA6 powder. The composite truss cylindrical lattice structures with different configurations are presented in Figure 8. Four types of structural configurations were investigated in this study, and each structure consisted of six layers. The geometric dimensions corresponding to different relative densities are listed in Table 1. The maximum dimensional error between the designed values and measured specimens was less than 0.4\%.

\noindent \includegraphics*[width=3.76in, height=3.15in]{image17}

\noindent Fig. 8  Photographs showcasing the composite truss cylindrical lattice structures

\noindent 
\paragraph{4.2 Material micro-morphology}

\noindent \includegraphics*[width=4.62in, height=3.02in]{image18}

\noindent Fig. 9  Microstructure of nylon PA6 with short glass fibers composite material

Figure 9 shows the micromorphology of short glass fiber reinforced nylon PA6 composite. It can be observed that the glass fibers present a slender filament shape with uniform diameter distribution and are well dispersed within the nylon PA6 matrix. Moreover, tight interfacial bonding between glass fibers and the nylon PA6 matrix is observed, with no obvious gaps or debonding detected. This indicates that load can be effectively transferred between the two phases.

\noindent 
\paragraph{4.3 Characterization of matrix materials}

To characterize the mechanical properties of short glass fiber reinforced nylon PA6, standard specimens were fabricated via the same printing method adopted for lattice samples with a printing direction of 90$\mathrm{{}^\circ}$. Tensile and compressive properties of the matrix material were tested in accordance with ASTM D638-14 and ASTM D695-15, respectively. The density and Poisson's ratio of the composite were 1.15 g/cm${}^{3}$ and 0.3. Tensile and compressive tests were carried out on the printed standard specimens using an INSTRON 5569 universal testing machine with a maximum load capacity of 50 kN at a nominal strain rate of 10${}^{-3}$s${}^{-1}$. The true stress-strain curves of tensile and compressive specimens are plotted in Figure 10, and the relevant parameters are summarized in Table 2.

\noindent \includegraphics*[width=2.84in, height=2.28in]{image19}\includegraphics*[width=2.81in, height=2.27in]{image20}

Fig. 10  Results of mechanical testing for the Nylon PA6 with glass fibers substrate: (a) Tensile stress-strain curve; (b) Compressive stress-strain curve.

Table 2 Mechanical properties of Nylon PA6 with short glass fibers substrate

\begin{tabular}{|p{1.3in}|p{1.3in}|p{1.3in}|} \hline 
Mechanical parameters & Unit symbol & Value \\ \hline 
Tensile modulus & MPa & 3944.36 \\ \hline 
Tensile strength & MPa & 70.47 \\ \hline 
Compression modulus & MPa & 3066.36 \\ \hline 
Compression strength & MPa & 82.47 \\ \hline 
\end{tabular}


\paragraph{4.4 Quasi-static compression test of composite truss cylindrical lattice structures}

To explore the mechanical behaviors of composite truss cylindrical lattice structures during the quasi-static elastic deformation stage, quasi-static compression tests were performed on four groups of lattice structures with different configurations using an Instron 5569 universal electronic testing machine. This apparatus is equipped with an electronic displacement controller, load sensor, camera, loading platform and two rigid compression platens. During the tests, each lattice specimen was placed between the upper and lower rigid platens. The upper platen was driven by the displacement controller and moved downward at a constant rate of 1 mm/min, while the lower platen remained fixed. As shown in Figure 11, the testing machine automatically collected and recorded compression load data throughout the tests, and a camera was used to capture the real-time deformation process of specimens. Combined with the original geometric dimensions of specimens, the stress and strain values under uniform compression were calculated. The elastic modulus of composite truss cylindrical lattice structures was determined from the slope of the stress-strain curve within the small elastic deformation region.

\noindent \includegraphics*[width=5.14in, height=3.33in]{image21}

Fig. 11  Schematic diagram of different composite truss cylindrical lattice structures under quasi-static compression experiments

\noindent 
\paragraph{4.5 Low-velocity impact test of composite truss cylindrical lattice structures}

Drop weight impact tests were conducted on composite truss cylindrical lattice structures in this section using an Instron Ceast 9450 electronic impact tester. The experimental setup is illustrated in Figure 12. This device mainly consists of a drop weight unit, operation control panel, high-speed camera and data acquisition system. The drop weight is composed of an impact head and attachable mass blocks, and is fitted with a load cell with a measuring range of 90 kN.

\noindent \includegraphics*[width=4.16in, height=3.15in]{image22}

\noindent Fig. 12  Schematic diagram of the impact experimental setup for composite truss cylindrical lattice structures

\noindent 
\subsection{5 Mechanical properties of composite truss cylindrical lattice structures}

\noindent 
\paragraph{5.1 Mechanical properties of composite truss cylindrical lattice structures under quasi-static compression}

The collected force and displacement data were processed to obtain the compressive stress-strain curves of composite truss cylindrical lattice structures, as shown in Figure 13. The mechanical response curves exhibit three distinct stages: linear elasticity, plastic yielding and buckling failure. The mechanical properties of the structures degrade as the unit cell size increases. When the diameter of truss rods remains constant, an increase in unit cell size reduces the relative density of the lattice structures, which further lowers their compressive stiffness and strength. Meanwhile, a larger unit cell leads to a higher slenderness ratio of the truss rods and a lower peak stress. This indicates that lattice structures with lower relative density are more prone to buckling failure.

The compressive modulus and compressive strength of truss cylindrical lattice structures increase significantly with the rise of relative density, demonstrating that relative density is a critical geometric parameter for regulating their load-bearing capacity. Further analysis reveals that lattice structures with different topological configurations present distinct mechanical properties at the same relative density. Specifically, the face-centered cubic cylindrical lattice structure retains high stiffness and strength even at low relative densities, while the tesseract lattice structure undergoes buckling and instability at an earlier stage. The findings provide valuable references for the design of lightweight load-bearing components adopting composite truss cylindrical lattice structures.

\noindent \includegraphics*[width=5.76in, height=1.83in]{image23}

\noindent \includegraphics*[width=5.66in, height=1.79in]{image24}

\noindent \includegraphics*[width=5.75in, height=1.83in]{image25}

\noindent \includegraphics*[width=5.76in, height=1.84in]{image26}

\noindent Fig. 13  The stress-strain curves of the composite truss cylindrical lattice structures under quasi-static compression at different relative densities: (a) body-centered; (b) face-centered; (c) octahedral; (d) tesseract.

The compressive mechanical response curves obtained from finite element simulations are presented in Figure 13(a)-(d). It can be seen that the simulated compressive modulus and strength show good agreement with the experimental results. Lattice structures of the same type but different sizes exhibit similar mechanical responses, which consist of a typical linear elastic stage, a peak stress after buckling deformation, and subsequent structural failure accompanied by a decline in load. This trend is consistent with the mechanical characteristics of lattice structures reported in previous studies. As the unit cell size decreases, the relative density increases, leading to higher load values on the response curves. Since lattice structures of the same type undergo similar deformation and failure modes under compression, only the deformation processes of truss cylindrical lattice structures with a unit cell height of 15 mm are displayed in Figure 14(a)-(d). The compressive buckling deformation of struts observed in numerical simulations matches the experimental observations. Nevertheless, no fracture criterion was defined for the base material in the finite element model, so no fracture of nodes or struts occurs in the simulation results.

For the body-centered cubic cylindrical lattice structure, buckling firstly occurs on the vertical struts while no deformation is observed on the inclined struts. As the compressive strain increases, the vertical struts continue to buckle and the connecting nodes deform, followed by deformation of the inclined struts. Correspondingly, the compressive stress rapidly reaches its peak value. Further loading causes fracture of both vertical and inclined struts, leading to a sharp drop in stress on the stress-strain curve. Figure 14(a) shows the deformation contour of strut buckling at a strain of $\varepsilonup$=0.02, which clearly illustrates the buckling of vertical struts. With ongoing compression, the vertical struts experience progressive buckling and mid-strut fracture until the entire body-centered cubic cylindrical lattice structure collapses completely. Similarly, the deformation of face-centered cubic and tesseract cylindrical lattice structures follows the same pattern: vertical struts buckle first, and inclined struts buckle subsequently. Differing from the above three configurations, the deformation process of octahedral cylindrical lattice structure is shown in Figure 14(c). Composed of octahedral truss units, this structure possesses a higher density of strut connections. At the initial loading stage, loads are uniformly distributed through combined tension and compression of multi-directional inclined struts. When the strain rises to 0.025, fracture only takes place in local struts. Such phenomenon is attributed to the unique geometric arrangement of struts in octahedral cylindrical lattice structures, which effectively prevents overall buckling deformation.

\noindent \includegraphics*[width=5.71in, height=2.15in]{image27}

\noindent \includegraphics*[width=5.71in, height=2.14in]{image28}\includegraphics*[width=5.69in, height=2.14in]{image29}\includegraphics*[width=5.70in, height=2.16in]{image30}

Fig. 14  Deformation process of four kinds of composite truss cylindrical lattice structures from axial crushing numerical simulation: (a) body-centered; (b) face-centered; (c) octahedral; (d) tesseract.

\noindent \includegraphics*[width=5.85in, height=1.83in]{image31}

\noindent \includegraphics*[width=5.87in, height=1.83in]{image32}

\noindent \includegraphics*[width=5.84in, height=1.82in]{image33}

\noindent \includegraphics*[width=5.88in, height=1.83in]{image34}

Fig. 15  Comparison of theoretical predictions, numerical simulations and experimental results of different relative density truss cylindrical lattice structures: (a) body-centered; (b) face-centered; (c) octahedral; (d) Tesseract.

The compressive modulus of composite truss cylindrical lattice structures was calculated from the linear elastic segment of experimental stress-strain curves, and the compressive strength was determined by the peak stress on the curves. In terms of load-bearing capacity, the body-centered cubic and tesseract cylindrical lattice structures exhibit the poorest performance at low relative densities, with average compressive moduli of 43.29 MPa and 42.89 MPa, and compressive strengths of 0.85 MPa and 0.87 MPa, respectively. The face-centered cubic cylindrical lattice structure delivers the optimal load-bearing performance, with a compressive modulus of 85.18 MPa and a compressive strength of 1.52 MPa. From the perspective of lightweight design, the body-centered cubic and tesseract structures contain fewer struts. Given the identical strut diameter, the tesseract cylindrical lattice structure features the lowest relative density. The relationships between mechanical properties and relative density of the truss cylindrical lattice structures are plotted in Figure 15(a)-(d). According to the Gibson-Ashby power-law relation between material performance and relative density, the mechanical properties of the four proposed structures show an approximately linear correlation with relative density, indicating that these structures are tension-dominated. This characteristic meets the original design objective of high load-bearing capacity. Taking the octahedral cylindrical lattice structure with a relative density of 16\% as an example, its compressive modulus and compressive strength reach 117.53 MPa and 2.53 MPa, respectively.

The simulated compressive modulus and compressive strength were derived from the curves obtained via finite element numerical simulation. The mechanical properties of lattice structures calculated by theoretical methods, finite element results and experimental data are compared in Table 3. The average errors of compressive modulus and compressive strength among theoretical predictions, finite element simulations and experimental tests are less than 15\%, which verifies the reliability of the theoretical model. Accordingly, the established finite element model is applicable to the simulation and analysis of the strength and stiffness of truss lattice structures. The discrepancies mainly stem from three aspects. First, geometric defects exist between the actual printed specimens and ideal design models, including minor variations in strut diameter and wall thickness. Manufacturing imperfections such as fluctuating strut diameter, powder accumulation at joints and uneven wall thickness during 3D printing contribute greatly to the deviations among experimental, theoretical and numerical results. Second, the theoretical model fails to take into account stress concentration at joints and initial imperfections of struts. Third, there are differences between the actual boundary conditions in experiments and the ideal constraints adopted in finite element models.

Table 3  Comparison of theoretical, finite element simulation and experimental results for quasi-static compression of composite truss cylindrical lattice structures

\begin{tabular}{|p{0.5in}|p{0.4in}|p{0.4in}|p{0.5in}|p{0.6in}|p{0.4in}|p{0.4in}|p{0.5in}|p{0.6in}|p{0.4in}|} \hline 
Cylindrical lattice structures & Relative density & \multicolumn{4}{|p{1.9in}|}{Compression modulus (MPa)} & \multicolumn{4}{|p{1.9in}|}{Compression strength (MPa)} \\ \hline 
 &  & Theory & Simulation & Experiment & Average error(\%) & Theory & Simulation & Experiment & Average error\newline (\%) \\ \hline 
Body-centered & 0.057 & 42.12 & 45.43 & 43.29\newline $\mathrm{\pm}$2.68 & 2.78 & 0.74 & 0.89 & 0.81\newline $\mathrm{\pm}$0.04 & 6.28 \\ \hline 
 & 0.092 & 72.97 & 76.23 & 78.25\newline $\mathrm{\pm}$2.16 & 2.50 & 1.81 & 1.65 & 1.73\newline $\mathrm{\pm}$0.02 & 3.08 \\ \hline 
 & 0.13 & 98.92 & 97.74 & 98.16\newline $\mathrm{\pm}$1.34 & 0.44 & 2.29 & 2.18 & 2.72\newline $\mathrm{\pm}$0.14 & 8.99 \\ \hline 
Face-centered & 0.072 & 86.56 & 84.72 & 85.18\newline $\mathrm{\pm}$2.84 & 0.84 & 1.58 & 1.54 & 1.52\newline $\mathrm{\pm}$0.11 & 1.44 \\ \hline 
 & 0.113 & 142.25 & 146.47 & 137.73\newline $\mathrm{\pm}$1.30 & 2.07 & 3.28 & 2.69 & 2.41\newline $\mathrm{\pm}$0.07 & 11.62 \\ \hline 
 & 0.156 & 171.07 & 170.42 & 166.74\newline $\mathrm{\pm}$3.39 & 1.05 & 4.78 & 4.07 & 4.15\newline $\mathrm{\pm}$0.10 & 6.87 \\ \hline 
Octahedral & 0.0999 & 76.39 & 70.51 & 70.66\newline $\mathrm{\pm}$1.32 & 3.56 & 1.81 & 1.7 & 1.58\newline $\mathrm{\pm}$0.06 & 4.58 \\ \hline 
 & 0.16 & 129.50 & 123.74 & 117.53\newline $\mathrm{\pm}$5.23 & 3.27 & 2.81 & 2.91 & 2.53\newline $\mathrm{\pm}$0.20 & 5.33 \\ \hline 
 & 0.229 & 140.37 & 144.65 & 138.54\newline $\mathrm{\pm}$3.30 & 1.64 & 3.47 & 3.61 & 3.31\newline $\mathrm{\pm}$0.33 & 2.95 \\ \hline 
Tesseract & 0.051 & 42.35 & 43.48 & 42.89\newline $\mathrm{\pm}$0.75 & 0.89 & 0.88 & 0.93 & 0.87\newline $\mathrm{\pm}$0.03 & 2.74 \\ \hline 
 & 0.073 & 73.69 & 78.48 & 76.47\newline $\mathrm{\pm}$2.96 & 2.21 & 1.83 & 2.06 & 2.09\newline $\mathrm{\pm}$0.14 & 5.46 \\ \hline 
 & 0.091 & 100.30 & 109.49 & 107.27\newline $\mathrm{\pm}$4.37 & 3.40 & 2.998 & 2.64 & 2.97\newline $\mathrm{\pm}$0.10 & 5.33 \\ \hline 
\end{tabular}


\paragraph{5.2 Energy absorption characteristics of composite truss cylindrical lattice structures under impact loading}

Figure 16 compares the numerical and experimental results of composite truss cylindrical lattice structures under low-velocity impact loading. The mechanical responses of four types of cylindrical lattice configurations during impact are illustrated in this figure. As listed in Table 4, the maximum average error of the peak load between finite element simulation and experimental tests is less than 5.61\%. Local offset in the load-displacement curves is mainly attributed to slight differences in initial contact conditions during testing, such as minor initial gaps between specimens and impact heads, as well as the simplification of strain rate sensitivity in the material constitutive model${}^{[51-52]}$. Overall, the established finite element model for low-velocity impact is capable of simulating the dynamic mechanical performance of truss cylindrical lattice structures. Benefiting from uniform strut distribution and dense connecting nodes, the face-centered cubic cylindrical lattice structure possesses superior stiffness and stability under impact, presenting a gently declining load curve. In contrast, the body-centered cubic and tesseract cylindrical lattice structures have fewer struts, which leads to earlier buckling and instability under the same impact energy, accompanied by a rapid drop in load. Owing to its high density of strut connections, the octahedral cylindrical lattice structure realizes coordinated loading through multi-directional inclined struts during impact, thus delaying the overall failure of the structure.

\noindent \includegraphics*[width=2.80in, height=2.38in]{image35}\includegraphics*[width=2.80in, height=2.35in]{image36}

\noindent \includegraphics*[width=2.85in, height=2.39in]{image37} \includegraphics*[width=2.75in, height=2.38in]{image38}

\noindent Fig. 16  Numerical simulation and experimental test comparison of composite truss cylindrical lattice structures under the low-speed impact: (a) body-centered; (b) face-centered; (c) octahedral; (d) tesseract.

To evaluate the energy absorption performance of composite cylindrical lattice structures under impact loading, the conventional energy absorption ($E_{a} $) index is not normalized by mass and thus cannot objectively reflect the intrinsic energy absorption capacity. Therefore, the specific energy absorption ($W_{SEA} $) is defined as the total absorbed energy divided by the structural mass, which is expressed as:
\begin{equation} \label{GrindEQ__31_} 
W_{SEA} =\frac{E_{a} }{m}  
\end{equation} 
where $m$ denotes the mass of the composite cylindrical lattice structure.

Table 4  Comparison of peak loads between the finite element model and the experiment of the truss cylindrical lattice structures obtained through low-speed impact

\begin{tabular}{|p{0.8in}|p{0.8in}|p{0.8in}|p{0.8in}|p{0.8in}|} \hline 
Cylindrical lattice structures & Relative density $\bar{\rho }$ & Experiment(kN) & Simulation(kN) & Average error(\%) \\ \hline 
Body-centered & 0.130 & 12.89 & 11.57 & 2.05 \\ \hline 
 &  & 11.22 &  &  \\ \hline 
Face-centered & 0.156 & 12.04 & 11.50 & 4.34 \\ \hline 
 &  & 13.20 &  &  \\ \hline 
Octahedral & 0.229 & 12.74 & 11.74 & 5.61 \\ \hline 
 &  & 13.53 &  &  \\ \hline 
Tesseract & 0.091 & 7.62 & 7.44 & 1.26 \\ \hline 
 &  & 7.64 &  &  \\ \hline 
\end{tabular}

Figure 17 compares the energy absorption performance of four types of composite truss cylindrical lattice structures under impact loading. The masses of the body-centered cubic, face-centered cubic, octahedral and tesseract cylindrical lattice structures are 39.1 g, 58.4 g, 65.7 g and 50.5 g, respectively. As indicated in the figure, the face-centered cubic cylindrical lattice structure achieves the highest total energy absorption and specific energy absorption. Despite its lower mass, the body-centered cubic cylindrical lattice structure exhibits higher energy absorption efficiency than the octahedral counterpart. This result further demonstrates the significance of specific energy absorption in evaluating the combined performance of lightweight design and energy absorption capacity for lattice structures.

\noindent \includegraphics*[width=2.77in, height=2.35in]{image39}\includegraphics*[width=2.81in, height=2.35in]{image40}

\noindent Fig. 17  Energy absorption values of different truss cylindrical lattice structures: (a) Energy absorption and (b) Specific energy absorption.

\noindent 
\paragraph{5.3 Comparison of mechanical properties with other common truss lattice structures}

Normalization analysis was conducted on the mechanical properties of lattice structures to eliminate the influence of matrix material parameters, and the results are presented in Figure 18. It is observed that among the four configurations, the face-centered cubic cylindrical lattice structure delivers the optimal mechanical performance, while the body-centered cubic counterpart exhibits relatively inferior properties. Furthermore, the four proposed cylindrical lattice structures were compared with those reported in existing literature. The present structures possess higher relative compressive modulus than conventional lattice configurations. In terms of relative compressive strength, the tesseract cylindrical lattice structure shows prominent load-bearing advantages at low relative densities. By contrast, the body-centered cubic and octahedral cylindrical lattice structures exhibit comparable load-bearing capacity to other lattice structures.

\noindent \includegraphics*[width=2.80in, height=2.42in]{image41}\includegraphics*[width=2.80in, height=2.39in]{image42}

\noindent \includegraphics*[width=3.55in, height=1.51in]{image43}

\noindent Fig. 18  Comparison of normalized mechanical properties with other lattice structures: (a) Relative compressive modulus; (b) Relative compressive strength.

\noindent 
\subsection{6 Conclusions}

This paper extends the design from traditional truss cubic lattice structures to truss cylindrical lattice structures, and establishes a theoretical model for composite truss cylindrical lattice structures at low relative densities. The accuracy of the theoretical model is verified through a combination of numerical simulation and experimental tests, and the mechanical properties and deformation characteristics of the structures are systematically investigated. In addition, the energy absorption performance of truss cylindrical lattice structures under impact loading is studied. The main conclusions are summarized as follows:

\eqref{GrindEQ__1_} Analytical expressions for the relative density are derived for truss cylindrical lattice structures with four typical topological configurations, namely body-centered cubic, face-centered cubic, octahedral and tesseract lattice structures. Based on the energy method, theoretical prediction models for the equivalent compressive modulus and compressive strength of different truss lattice structures are established.

\eqref{GrindEQ__2_} The effects of relative density and topological configuration on mechanical properties are systematically revealed. The results show that the compressive modulus and compressive strength of the four types of truss cylindrical lattice structures under quasi-static compression are linearly and positively correlated with relative density. At low relative densities (less than 10\%), face-centered cubic and tesseract cylindrical lattice structures exhibit outstanding load-bearing capacity, with their compressive modulus and strength remarkably higher than those of body-centered cubic counterparts at the same density.

\eqref{GrindEQ__3_} Under drop-weight impact loading, the specific energy absorption of composite face-centered cubic cylindrical lattice structures is 1.33 to 3.56 times that of the other truss cylindrical lattice structures, demonstrating its excellent impact resistance.

\eqref{GrindEQ__4_} As illustrated in the Ashby material selection chart, where the horizontal and vertical axes represent the relative density and relative compressive modulus and strength of truss cylindrical lattice structures respectively, the proposed structures outperform many conventional lattice structures.

\noindent \textbf{Acknowledgments}

The present work was supported by the National Natural Science Foundation of China (NSFC) under [grant numbers: 12472345]. JX was support by the Scientific Research Inno-vation Capability Support Project for Young Faculty [grant numbers: ZYGXQNJSKYCXNLZCXM-M9]. The Open Project Program of Ministry of Education Key Laboratory for Advanced Textile Composite Materials [grant numbers: MATC 2024-Z02].

\noindent 
\subsection{References}

\noindent [1] Ali M, Sajjad U, Hussain I, et al. On the assessment of the mechanical properties of additively manufactured lattice structures[J]. Engineering Analysis with Boundary Elements, 2022, 142: 93-116.

\noindent [2] Yin H F, Zhang W Z, Zhu L C, et al. Review on lattice structures for energy absorption properties[J]. Composite Structures, 2023, 304: 116397.

\noindent [3] Wang P, Guo J Z, Yuan Y H, et al. Bio-inspired vertex-offset lattice metamaterials with enhanced stress stability and energy absorption[J]. Thin-Walled Structures, 2025, 210: 113060.

\noindent [4] Tian Y Y, Zhang X, Jarl\"{o}v A, et al. Programmable heterogeneous lamellar lattice architecture for dual mechanical protection[J]. Proceedings of the National Academy of Sciences. 2024, 121: e2407362121.

\noindent [5] Khan N, Riccio A. A systematic review of design for additive manufacturing of aerospace lattice structures: Current trends and future directions. Progress in Aerospace Sciences. 2024, 149: 101021.

\noindent [6] Jiao P C. Mechanical energy metamaterials in interstellar travel[J]. Progress in Materials Science, 2023, 137: 101132.

\noindent [7] Bai J B, Li S L, Liu T W, et al. Variable camber structure for aircraft wing based on zero Poisson's ratio lattice structure[J]. Journal of the Mechanics and Physics of Solids, 2025, 203, 106241.

\noindent [8] Blakey-Milner B, Gradl P, Snedden G, et al. Metal additive manufacturing in aerospace: A review[J]. Material \& Design, 2021, 209, 110008.

\noindent [9] Caton M, Kowalewski J F, Guo J N, et al. Bridging hard and soft: Mechanical metamaterials enable rigid torque transmission in soft robots[J]. Science Robotics, 2025, 10: eads0548.

\noindent [10] Li Q, Wang W, Tan H, et al. Energy absorption characteristics of modular assembly structures under quasi-static compression load[J]. Composite Structures, 2024, 342: 118260.

\noindent [11] Ma W W S M, Yang H, Zhao Y J, et al. Multi-physical lattice metamaterials enabled by additive manufacturing: design principles, interaction mechanisms, and multifunctional applications[J]. Advanced Science, 2025, 12: 2405835.

\noindent [12] Meng Z Z, Xu Y D, Zhou W, et al. Foamed cementitious composites with 3D-printed auxetic lattice reinforcement: enhancing static and cyclic performance[J]. Composites Part B: Engineering, 2025, 303: 112614.

\noindent [13] Garrido C, Pincheira G, Valle R, et al. Plastic deformation behavior and energy absorption performance of a composite metamaterial based on asymmetric auxetic lattices[J]. Composite Structures, 2024, 346: 118410.

\noindent [14] Zhang Q, Li F Y, Zhu D C, et al. Quasi-static and impact performance study of a three-dimensional negative Poisson's ratio structure with adjustable mechanical properties[J]. International Journal of Impact Engineering, 2024, 193: 105057.

\noindent [15] Chen X Y, Moughames J, Ji Q X, et al. 3D lightweight mechanical metamaterial with nearly isotropic inelastic large deformation response[J]. Journal of the Mechanics and Physics of Solids, 2022, 169: 105057.

\noindent [16] Li X W, Tan Y H, Wang P, et al. Metallic microlattice and epoxy interpenetrating phase composites: Experimental and simulation studies on superior mechanical properties and their mechanisms[J]. Composites Part A: Applied Science and Manufacturing, 2020, 135: 105934.

\noindent [17] White B C, Garland A P, Alberdi R, et al. Interpenetrating lattices with enhanced mechanical functionality[J]. Additive Manufacturing, 2021, 38: 101741.

\noindent [18] Zhao M, Cui J M, Chen L, et al. Enhanced mechanical properties and energy absorption of lattice metamaterials inspired by crystal imperfections[J]. Composite Structures, 2025, 356: 118894.

\noindent [19] Guo X, Wang E D, Yang H, et al. Mechanical characterization and constitutive modeling of additively-manufactured polymeric materials and lattice structures[J]. Journal of the Mechanics and Physics of Solids, 2024, 189: 105711.

\noindent [20] Tancogne-Dejean T, Mohr D. Elastically-isotropic truss lattice materials of reduced plastic anisotropy[J]. International Journal of Solids and Structures, 2018, 138: 24-39.

\noindent [21] Chen Z, Mordehai D. Homogenization of the yield surface of body-centered cubic and octet lattice structures under multiaxial loadings[J]. International Journal of Solids and Structures, 2025, 320: 113486.

\noindent [22] Kaya G, Y$\mathrm{\imath}$ld$\mathrm{\imath}$z F, Solak K, et al. An experimental and FEA investigation of deformation characteristics of additively manufactured Ti6Al4V lattice structures. European Journal of Mechanics - A/Solids[J], 2025, 112: 105657.

\noindent [23] Yuan Y, Liu X N, Zhang Y F, et al. Modified face-centred cubic lattice with enhanced mechanical properties[J]. Thin-walled structures, 2024, 315: 118440.

\noindent [24] Zhang S K, Yu K H, Dong L. Compressive property of a hybrid hierarchical metamaterial[J]. Materials Today Communications, 2022, 33: 104260.

\noindent [25] Sun Z P, Zhang J J, Hua T, et al. Experimental study on the post-yield responses and crushing patterns of additively manufactured Rhombic Dodecahedron lattices in biaxial compression[J]. Thin-walled structures, 2025, 213: 113331.

\noindent [26] Fallahi H, Sadighi M, Samandari S S. Compressive mechanical properties of open-cell cubic and rhombic dodecahedron lattice structures of variable cross section fabricated by fused deposition modeling[J]. Mechanics of Materials, 2023, 184: 104712.

\noindent [27] Rodrigo C, Xu S Q, Durandet Y, et al. Mechanical response of functionally graded lattices with different density grading strategies[J]. Thin-walled structures, 2023, 192: 111132.

\noindent [28] Cao X F, Duan S Y, Liang J, et al. Mechanical properties of an improved 3D-printed rhombic dodecahedron stainless steel lattice structure of variable cross section[J]. International Journal of Mechanical Sciences, 2018, 145: 53-63.

\noindent [29] Xiao L J, Xu X, Feng G Z, et al. Compressive performance and energy absorption of additively manufactured metallic hybrid lattice structures[J]. International Journal of Mechanical Sciences, 2022, 219: 107093.

\noindent [30] Li J W, Bai Y, Liu E D, et al. Investigation of mechanical properties for bio-inspired lattice metamaterials with crossed struts in multiple units[J]. Materials Today Communications, 2025, 47: 112963.

\noindent [31] Chen L, Chen T, Feng S, et al. Self-similar nesting strategy enables lattices achieve dual energy-absorbing plateaus[J]. International Journal of Mechanical Sciences, 2024, 276: 109445.

\noindent [32] Han T, Qi D D, Ma J, et al. Generative design and mechanical properties of the lattice structures for tensile and compressive loading conditions fabricated by selective laser melting[J]. Mechanics of Materials, 2024, 188: 104840.

\noindent [33] Kothari T, Liu X, Kobir M H, et al. Design and analysis of reinforced I-beam body-centered cubic lattices for enhanced energy absorption[J]. Composite Structures, 2025, 188: 104840.

\noindent [34] Sasagawa T, Nimura N, Tanaka M. Flexible lattice structure using curved struts based on body-centered cubic structure[J]. Additive Manufacturing, 2025, 103: 104746.

\noindent [35] Chen Z, Sahoo S, P\'{e}rez-Prado M T, et al. The extended scaling laws of the mechanical properties of additively manufactured body-centered cubic lattice structures under large compressive strains[J]. Mechanics of Materials, 2024, 196: 105075.

\noindent [36] Zhang P J, Chen X Y, Yu P H, et al. Grid hollow octet truss lattices that are stable at low relative density[J]. Journal of the Mechanics and Physics of Solids, 2025, 197: 106068.

\noindent [37] Lan T, Tran P, Nguyen K T, H. et al. Improving energy absorption and failure characteristic of additively manufactured lattice structures using hollow and curving techniques[J]. Composite Structures, 2025, 188: 104840.

\noindent [38] Wang P, Yang F, Li P F, et al. Bio-inspired vertex modified lattice with enhanced mechanical properties[J]. International Journal of Mechanical Sciences, 2023, 244: 108081.

\noindent [39] Wang P, Yuan Y H, Guo J Z, et al. Bio-inspired vertex-offset lattices with high stress and stable deformation under axial and oblique impacts[J]. International Journal of Impact Engineering, 2026, 207: 105484.

\noindent [40] Tancogne-Dejean T, Spierings A B, Mohr D. Additively-manufactured metallic micro-lattice materials for high specific energy absorption under static and dynamic loading[J]. Acta Materialia, 2016, 116: 14-28.

\noindent [41] Sun Z P, Guo Y B, Shim V P W. Static and dynamic crushing of polymeric lattices fabricated by fused deposition modelling and selective laser sintering-an experimental investigation[J]. International Journal of Impact Engineering, 2022, 160: 104059.

\noindent [42] Chen C, Wang H R, Jing Y J, et al. Load dispersion mechanism of a novel structure composed by pre-stretched fabrics and lattice[J]. International Journal of Impact Engineering, 2025, 206: 105455.

\noindent [43] Zimmerman B K, Grenley S P, Saunders A M, et al. Solid face sheets enable lattice metamaterials to withstand high-amplitude impulsive loading without yielding[J]. International Journal of Impact Engineering, 2025, 195: 105130.

\noindent [44] Gao Q, Liao W H, Wang L M. An analytical model of cylindrical double-arrowed honeycomb with negative Poisson's ratio[J]. International Journal of Mechanical Sciences, 2020, 173: 105400.

\noindent [45] Zhang L Y, Lei H S, Wang F, et al. Twist design of lattice structure fabricated by powder bed fusion to adjust the energy absorption behavior[J]. Composite Structures, 2024, 342: 118224.

\noindent [46] Tian Y, Wang H T, Chen Z, et al. Design principles for dual-phase lattice cylindrical tubes with excellent energy absorption capability[J]. Composite Structures, 2025, 359: 119015.

\noindent [47] Pan Y, Zhou Y, Wang M, et al. A novel reinforced cylindrical negative stiffness metamaterial for shock isolation: Analysis and application[J]. International Journal of Solids and Structures, 2023, 279: 112391.

\noindent [48] Li X Y, Xiao L J, Song W D. Deformation and failure modes of Ti-6Al-4V lattice-walled tubes under uniaxial compression[J]. International Journal of Impact Engineering, 2019, 130: 27-40.

\noindent [49] Sood M, Wu C M, Yang Y C. Mechanical properties of 3D-printed lattice cylindrical ptructure with recyclable elastomeric and thermoplastic polymers[J]. Joumal of Polymers and the Environment, 2024, 32: 4196-4212.

\noindent [50] Mars J., Chebbi E, Wali M, Dammak F. Numerical and experimental investigations of low velocity impact on glass fiber-reinforced polyamide[J]. Composites Part B, 2018, 146: 116-123.

\noindent [51] Iranmanesh N, Yazdani Sarvestani H, Ashrafi B, et al. Beyond honeycombs: Core topology's role in 3D-printed sandwich panels[J]. Materials Today Communications, 2023, 37: 107548.

\noindent [52] Chao X J, Hu H T, Lin J J, et al. Bioinspired twist-hyperbolic metamaterial for impact buffering and self-powered real-time sensing in UAVs[J]. Science Advances, 2025, 11: eadw6179.

\noindent [53] Ahmadi S M, Yavari S A, Wauthle R, et al. Additively manufactured open-cell porous biomaterials made from six different space-filling unit cells: the mechanical and morphological properties[J]. Materials, 2015, 8: 1871-1896.

\noindent [54] Kadkhodapour J, Montazerian H., Darabi A C, et al. Failure mechanisms of additively manufactured porous biomaterials: effects of porosity and type of unit cell[J]. Journal of the mechanical behavior of biomedical materials, 2015, 50: 180-191.

\noindent [55] Bai L, Gong C, Chen X H, et al. Heterogeneous compressive responses of additively manufactured Ti-6Al-4V lattice structures by varying geometric parameters of cells[J]. International Journal of Mechanical Sciences, 2022, 214: 106992.

\noindent [56] Leary M, Mazur M, Williams H, et al. Inconel 625 lattice structures manufactured by selective laser melting (SLM): mechanical properties, deformation and failure modes[J]. Materials \& Design, 2018, 157: 179-199.


\end{document}

